% ============================================================
% Cheat sheet, argomento: il pattern Observer (dispensa, capitolo 15)
% ============================================================
\argomento{Observer}{Dispensa: capitolo 15. Pattern comportamentale della GoF}[Pattern Observer]

\scheda{Carta d'identità}{
  \pintento{Definire una dipendenza \textbf{1:N} fra oggetti: quando il \textbf{soggetto} cambia stato, tutti gli
    \textbf{osservatori} vengono avvisati e aggiornati \textbf{automaticamente}.}
  \ptesto{«quando X cambia, Y deve aggiornarsi automaticamente», «notificare», «avvisare tutti gli interessati»,
    «iscriversi» e «cancellarsi», «abbonati», «più viste degli stessi dati», «osservatori di tipi diversi, aggiunti
    mentre il programma gira», «un sensore o un servizio fornisce un dato aggiornato», «eventi dell'interfaccia
    grafica»; nelle slide: aggiornare il valore delle azioni quando cambia la quotazione della società.}
  \pcodice{il soggetto ha una \textbf{lista di un'interfaccia} (\code{List<Osservatore>}), un metodo per
    \textbf{registrare} e uno per \textbf{togliere} gli osservatori (\code{attach}/\code{detach},
    \code{add\ldots Listener}), e un metodo di \textbf{notifica} che percorre la lista chiamando \code{update()} su
    tutti; il metodo che cambia lo stato (\code{setStato}) chiama la notifica. L'osservatore implementa l'interfaccia,
    ha un riferimento al soggetto e nel suo \code{update()} ne legge il nuovo stato.}
  \pdiagramma{\code{Subject} $\to$ \code{«interface» Observer} con molteplicità \code{*} (ruolo \code{observers});
    \code{ConcreteObserver} realizza Observer e ha un'associazione \code{1} verso \code{ConcreteSubject}; spesso la
    nota \code{for all o in observers: o.update()}.}
  \pvarianti{soggetto astratto (classe con la lista) o interfaccia; i \textit{listener} di Swing con una
    \textit{lambda} (\code{pulsante.addActionListener(e -> \ldots)}); il modello \textit{pull} (l'osservatore chiede lo
    stato, come qui) o \textit{push} (il soggetto passa il dato a \code{update(valore)}) \oltre; il cuore dell'MVC (la
    View osserva il Model). \code{java.util.Observable} esiste ma è deprecata da Java 9. In Java il metodo non può
    chiamarsi \code{notify()}: esiste già, \code{final}, in \code{Object}.}
  \pconfondere{il \textbf{Composite}: anche il composto ha una lista e chiama un metodo su tutti, ma i figli hanno la
    \textbf{stessa interfaccia} del composto e formano un albero; lo \textbf{Strategy}: il contesto ha \textbf{un} oggetto
    di un'interfaccia, non una lista, e gli delega un calcolo.}
}

\begin{struttura}{observer-struttura}{La struttura e il diagramma di sequenza}
\begin{tikzpicture}
  \umlclass[3.0cm]{as}{(0,0)}{\textit{Subject}}{}{+ attach(o: Observer)\\+ detach(o: Observer)\\+ notify()}
  \umlclass[2.4cm]{ob}{(5.4,0)}{\stereotipo{interface}{Observer}}{}{+ update()}
  \umlclass[3.0cm]{cs}{(0,-2.8)}{ConcreteSubject}{- state}{+ getState()\\+ setState(s)}
  \umlclass[2.8cm]{co}{(5.4,-2.8)}{ConcreteObserver}{- observerState}{+ update()}
  \draw[navig] (as.east) -- (as.east -| ob.west);
  \estO{(as.east -| ob.west)}{*}{observers}
  \draw[gen] (cs.north) -- (as.south);
  \draw[realizza] (co.north) -- (ob.south);
  \draw[navig] (co.west) -- (co.west -| cs.east);
  \estE{(co.west -| cs.east)}{1}{subject}
\end{tikzpicture}

\medskip
\begin{tikzpicture}
  \foreach \n/\x/\t in {q1/0/{:ConcreteObserver}, q2/3.6/{:ConcreteSubject}, q3/7.0/{sorgente}}{
    \node[partec, minimum width=2.4cm] (\n) at (\x,0) {\t};
    \draw[vita] (\n.south) -- ++(0,-3.7);}
  \draw[sinc] (0.08,-0.75) -- node[msgl]{attach(this)} (3.52,-0.75);
  \attiv{3.6}{-0.75}{-1.0}
  \draw[sinc] (6.92,-1.4) -- node[msgl]{setState(s)} (3.68,-1.4);
  \attiv{3.6}{-1.4}{-3.4}
  \draw[sinc] (3.68,-1.75) -- (4.6,-1.75) -- (4.6,-2.0) -- (3.68,-2.0);
  \node[font=\tiny, anchor=west] at (4.65,-1.87) {notify()};
  \draw[sinc] (3.52,-2.35) -- node[msgl]{update()} (0.08,-2.35);
  \attiv{0}{-2.35}{-3.3}
  \draw[sinc] (0.08,-2.7) -- node[msgl]{getState()} (3.52,-2.7);
  \draw[ritorno] (3.52,-3.05) -- node[msgl]{s} (0.08,-3.05);
\end{tikzpicture}
\tcblower
\textbf{Subject}: tiene la lista degli osservatori e li avvisa (\code{notify}). \textbf{Observer}: l'interfaccia con
\code{update()}. \textbf{ConcreteSubject}: ha lo stato; quando cambia, chiama \code{notify}. \textbf{ConcreteObserver}:
conosce il soggetto e, ricevuto \code{update()}, gli chiede il nuovo stato.

Sotto, la sequenza: l'osservatore si registra; più tardi una sorgente esterna cambia lo stato; il soggetto notifica
tutti, e ogni osservatore rilegge lo stato. Il soggetto conosce gli osservatori \textbf{solo} attraverso
l'interfaccia: se ne possono aggiungere di nuovi senza toccarlo.
\end{struttura}

\begin{programma}{observer-esempio}{Il problema delle slide: il valore delle azioni}
\begin{lstlisting}[style=javacs, name=Portafoglio.java]
public class Portafoglio {
    public static void main(String[] args) {
        Societa acme = new Societa("ACME", 10.0);
        AzioneAcquistata mie =
                new AzioneAcquistata(acme, 100);
        AzioneAcquistata tue =
                new AzioneAcquistata(acme, 40);
        System.out.println(mie + "\n" + tue);
        acme.setQuotazione(12.5); // la borsa cambia
        System.out.println(mie + "\n" + tue);
    }
}
interface OsservatoreQuotazione {       // Observer
    void update();
}
abstract class AbstractSocieta {        // Subject
    private final List<OsservatoreQuotazione>
            observers = new ArrayList<>();
    public void attach(OsservatoreQuotazione o) {
        observers.add(o);
    }
    public void detach(OsservatoreQuotazione o) {
        observers.remove(o);
    }
    protected void notifica() {         // il notify
        for (OsservatoreQuotazione o : observers)
            o.update();
    }
}
class Societa extends AbstractSocieta {
    private final String nome;
    private double quotazione;
    Societa(String nome, double quotazione) {
        this.nome = nome;
        this.quotazione = quotazione;
    }
    String getNome() { return nome; }
    double getQuotazione() { return quotazione; }
    void setQuotazione(double q) {
        quotazione = q;
        notifica();     // avvisa gli osservatori
    }
}
class AzioneAcquistata
        implements OsservatoreQuotazione {
    private final Societa societa;
    private final int quantita;
    private double valore;
    AzioneAcquistata(Societa s, int quantita) {
        societa = s;
        this.quantita = quantita;
        s.attach(this);           // si iscrive
        update();                 // valore iniziale
    }
    public void update() {
        valore = societa.getQuotazione() * quantita;
    }
    public String toString() {
        return quantita + " azioni "
                + societa.getNome() + " = " + valore
                + " euro";
    }
}
\end{lstlisting}
\tcblower
\begin{tikzpicture}
  \umlclass[3.4cm]{as}{(0,0)}{\textit{AbstractSocieta}}{}{+ attach(o)\\+ detach(o)\\\# notifica()}
  \umlclass[3.4cm]{oq}{(4.6,0)}{\stereotipo{interface}{OsservatoreQuotazione}}{}{+ update()}
  \umlclass[3.4cm]{so}{(0,-3.0)}{Societa}{- nome: String \{readOnly\}\\- quotazione: double}{\textasciitilde\ getQuotazione(): double\\\textasciitilde\ setQuotazione(q: double)}
  \umlclass[3.4cm]{az}{(4.6,-3.0)}{AzioneAcquistata}{- quantita: int \{readOnly\}\\- valore: double}{+ update()}
  \draw[navig] (as.east) -- (as.east -| oq.west);
  \estO{(as.east -| oq.west)}{*}{observers}
  \draw[gen] (so.north) -- (as.south);
  \draw[realizza] (az.north) -- (oq.south);
  \draw[navig] (az.west) -- (az.west -| so.east);
  \estE{(az.west -| so.east)}{1}{societa}
\end{tikzpicture}

\medskip
\begin{lstlisting}[style=shellcs, name=Terminale]
$ javac Portafoglio.java
$ java Portafoglio
100 azioni ACME = 1000.0 euro
40 azioni ACME = 400.0 euro
100 azioni ACME = 1250.0 euro
40 azioni ACME = 500.0 euro
\end{lstlisting}
\regola{Il \code{main} non chiede mai alle azioni di ricalcolarsi: cambia la quotazione, e le due azioni lo vengono a
sapere da sole. \code{AzioneAcquistata} si iscrive nel costruttore (\code{attach}); \code{setQuotazione} è un setter
con una riga in più, la notifica. Le slide danno il diagramma; il codice è della dispensa. L'esempio di codice
delle slide ha la stessa forma, con \code{Soggetto}, \code{Osservatore} e \code{aggiorna()}.}
\end{programma}
